% Complete integration and expectation proofs for P15-U1.
\section{Integration and expectation: proofs}
\label{sec:annex-p15-integration-expectation}

Use the definitions of Chapter~\ref{sec:p15-integration-expectation}.
The measure and simple-approximation prerequisites are in
Chapters~\ref{sec:p15-measurable-probability}
and~\ref{sec:p15-measure-extension}. Nonnegative sums and products
use their extended conventions throughout
\cite[secs.~1.4--1.6]{DurrettPTEAuthor2019}.

\begin{annexitem}{annex:p15.simple-integral-laws}{thm:p15-simple-integral-laws}{Refining simple functions}
\begin{proof}
Suppose \(s\) is represented on two finite disjoint measurable decompositions
\((A_i)\) and \((B_j)\), with respective values \(a_i,b_j\).
The sets \(C_{ij}=A_i\cap B_j\) form a common partition.
On a nonempty \(C_{ij}\), equality of the two functions gives
\(a_i=b_j\). Finite additivity of \(\mu\) and multiplication
by a fixed finite nonnegative scalar give
\[
 \sum_i a_i\mu(A_i)=\sum_{i,j}a_i\mu(C_{ij})
                 =\sum_{i,j}b_j\mu(C_{ij})=\sum_j b_j\mu(B_j).
\]
Empty intersections have measure zero. These finite distributive
identities remain valid when a measure is infinite: a zero coefficient
contributes zero on every piece, and a positive coefficient distributes
over nonnegative extended sums. This proves independence of representation.

Refine the partitions for any two nonnegative simple functions \(s,t\).
On each common cell their sum has the sum of their constant values,
and a scalar multiple has the corresponding multiple. Finite
distributivity therefore proves additivity and scaling of their integrals,
including the zero scalar. If \(s\leq t\), the values are ordered
on every nonempty cell; multiplication by its nonnegative measure
preserves that order. This proves monotonicity. The indicator formula
follows from the two parts \(A,A^c\), with respective coefficients
one and zero.

For a finite sum over possibly overlapping \(B_j\), partition
\(\Omega\) according to membership and nonmembership in every
\(B_j\). There are finitely many such measurable cells. On each
cell the value of the sum is the sum of exactly those coefficients
whose sets contain that cell. Distributing and regrouping the resulting
finite nonnegative sums gives \(\sum_j a_j\mu(B_j)\). This also
handles an empty original sum, which is the zero function.
\end{proof}
\end{annexitem}

\begin{annexitem}{annex:p15.monotone-integral}{thm:p15-monotone-integral}{The simple supremum and monotone convergence}
\begin{proof}
If \(f\) is nonnegative simple, every simple \(s\leq f\) has
integral at most its simple integral by the preceding theorem, and
\(s=f\) is itself admissible. Thus the two definitions agree.
If \(f\leq g\), each lower simple approximation to \(f\) is
one for \(g\); taking suprema proves monotonicity.

Let \(0\leq f_n\uparrow f\). Measurability of \(f\) follows
from the measurable-limit theorem. The integrals increase and are
bounded above by \(\int f\,d\mu\); let their extended limit be
\(L\). Fix a nonnegative simple \(s\leq f\), represented with
finite values \(a_j\) on a finite measurable partition \((A_j)\),
and choose \(0<\theta<1\). Put
\(E_n=\{f_n\geq\theta s\}\). These sets are measurable: on
each \(A_j\) this is the measurable level condition
\(f_n\geq\theta a_j\). They increase to \(\Omega\).
Where \(s=0\) the condition always holds; where \(s>0\), the
limit \(f\geq s>\theta s\) forces eventual membership.
The nonnegative simple function \(\theta s\mathbf1_{E_n}\)
is below \(f_n\), hence
\[
 \int f_n\,d\mu\geq\theta\sum_j a_j\mu(A_j\cap E_n).
\]
Continuity from below of the measure gives
\(\mu(A_j\cap E_n)\uparrow\mu(A_j)\). There are only finitely
many terms. A zero coefficient stays zero even if the measure tends
to infinity; a positive coefficient commutes with the increasing limit.
Consequently \(L\geq\theta\int s\,d\mu\).
If \(\int s\,d\mu=\infty\), this already forces \(L=\infty\).
If it is finite, let \(\theta\uparrow1\) to obtain
\(L\geq\int s\,d\mu\). Supremizing over all such \(s\)
gives \(L\geq\int f\,d\mu\). The earlier upper bound proves
equality, including the infinite case.
\end{proof}
\end{annexitem}

\begin{annexitem}{annex:p15.nonnegative-integral-laws}{thm:p15-nonnegative-integral-laws}{Nonnegative operations, exceptional sets and lower limits}
\begin{proof}
Choose nonnegative simple \(s_n\uparrow f\) and \(t_n\uparrow g\)
using Theorem~\ref{thm:p15-measurable-limits}. Their sums are simple
and increase to \(f+g\), with nonnegative extended addition.
This also proves measurability of that sum. Monotone convergence and
simple additivity give
\[
 \int(f+g)\,d\mu
   =\lim_n\left(\int s_n\,d\mu+\int t_n\,d\mu\right)
   =\int f\,d\mu+\int g\,d\mu.
\]
Increasing nonnegative sums commute with these limits, even when one
limit is infinite. For \(c>0\), use \(cs_n\uparrow cf\) to
obtain the scaling formula. For \(c=0\), \(cf\) is the zero
function by the adopted convention and both sides are zero.
Induction gives finite additivity of integrals, and the increasing
partial sums of \(\sum_n f_n\) then give the series formula by
monotone convergence.

If \(\int f\,d\mu=0\), the lower bound
\(f\geq k^{-1}\mathbf1_{\{f\geq1/k\}}\) gives
\(\mu(\{f\geq1/k\})=0\) for every positive integer \(k\).
Their union is \(\{f>0\}\), so it is null. Conversely, if
\(f=0\) outside a measurable null set \(N\), every nonnegative
simple \(s\leq f\) has each positive-value cell contained in
\(N\). Its integral is zero. The defining supremum is therefore
zero. If \(\int f\,d\mu=C<\infty\), the bound
\(f\geq n\mathbf1_{\{f=\infty\}}\) gives
\(\mu(\{f=\infty\})\leq C/n\) for every \(n\), so that set
is null.

Every nonnegative measurable function supported on a measurable null
set has zero integral by the result just proved. If \(f\leq g\)
outside such a set \(N\), split both functions into their parts on
\(N\) and \(N^c\). Additivity removes the zero-integral parts on
\(N\), and pointwise monotonicity compares the remaining parts.
This proves a.e. monotonicity and, by applying it twice, a.e. invariance.
For the a.e. version of monotone convergence, choose one measurable
null set outside which the asserted inequalities and convergence all
hold. If separate null sets were specified for each inequality, their
countable union is null. Set every function to zero on that union.
These replacements are measurable and have the same integrals; the
pointwise theorem applies to them.

For Fatou, put \(g_n=\inf_{m\geq n}f_m\). The functions are
nonnegative and measurable, increase to \(\liminf_n f_n\), and
satisfy \(\int g_n\,d\mu\leq\inf_{m\geq n}\int f_m\,d\mu\)
by monotonicity. Monotone convergence followed by the supremum over
\(n\) gives Fatou's inequality in the extended order.

Finally, \(\nu_f(\varnothing)=0\). If \(A_n\) are disjoint,
then pointwise
\(f\mathbf1_{\bigcup_n A_n}=\sum_n f\mathbf1_{A_n}\).
At most one term is nonzero at each point, including where \(f\)
is infinite. The series formula gives countable additivity of
\(\nu_f\), proving it is a measure on \(\mathcal F\).
\end{proof}
\end{annexitem}

\begin{annexitem}{annex:p15.signed-integral-laws}{thm:p15-signed-integral-laws}{Signed linearity and absolute control}
\begin{proof}
The parts of a measurable extended real function are measurable.
For \(a\geq0\), their sublevel sets are
\(\{f^+\leq a\}=\{f\leq a\}\) and
\(\{f^-\leq a\}=\{f\geq-a\}\); for \(a<0\) they are
empty. These are measurable by the level-set criterion and its strict
sublevel complements. The absolute value is the measurable nonnegative
sum of the parts. Thus all integrals in the definition have the stated
domains. Their sum equals \(\int|f|\,d\mu\), so integrability
is equivalent to finiteness of both part integrals. A.e. replacement
preserves the absolute value and both parts a.e., hence preserves their
integrals and integrability by the nonnegative theorem. In particular
the indicated finite representative has the same integral.

For nonnegative real integrable \(u,v\), the real function
\(f=u-v\) satisfies \(|f|\leq u+v\) and is integrable.
Pointwise \(u+f^-=v+f^+\). Nonnegative additivity gives
\[
 \int u\,d\mu+\int f^-\,d\mu
       =\int v\,d\mu+\int f^+\,d\mu.
\]
Every number here is finite, so rearrangement proves the asserted
difference-of-integrals formula.

For real integrable \(f,g\),
\(|f+g|\leq|f|+|g|\) proves integrability of the sum. Express it
as \((f^++g^+)-(f^-+g^-)\), a difference of two nonnegative
integrable real functions, and use the preceding formula and
nonnegative additivity to obtain signed additivity. For a positive
scalar the positive and negative parts scale by that scalar; for
a negative scalar they exchange roles and scale by its absolute value.
The zero scalar gives the zero function. This proves scaling for every
real scalar and hence the stated two-term linearity.

If an integrable real \(h\geq0\) a.e., its negative part is zero
a.e., so \(\int h\,d\mu\geq0\). Apply this to \(h=g-f\)
and use linearity to prove a.e. monotonicity. Since
\(-|f|\leq f\leq|f|\), monotonicity and scaling by minus one give
\(-\int|f|\,d\mu\leq\int f\,d\mu\leq\int|f|\,d\mu\),
which is the absolute-value bound. These facts also give the triangle
inequality for \(\|\cdot\|_1\), and its zero-class assertion is
the nonnegative zero-integral criterion applied to \(|f|\).

For disjoint measurable \(A_n\), apply nonnegative countable
additivity over sets to \(f^+\), \(f^-\) and \(|f|\). The
first two resulting series have finite sums bounded by the corresponding
part integrals on \(\Omega\). Therefore termwise subtraction is
legitimate and gives the asserted signed identity. Its absolute
convergence follows from
\[
 \sum_n\left|\int_{A_n}f\,d\mu\right|
       \leq\sum_n\int_{A_n}|f|\,d\mu
       =\int_{\bigcup_n A_n}|f|\,d\mu
       \leq\int|f|\,d\mu<\infty.
\]
\end{proof}
\end{annexitem}

\begin{annexitem}{annex:p15.integral-source-bridge}{thm:p15-integral-source-bridge}{Agreement with finite-support and sigma-finite constructions}
\begin{proof}
The hypotheses on \(f\) give
\(\int|f|\,d\mu\leq M\mu(E)<\infty\), so its integral is
finite. Every real simple function vanishing outside \(E\) is also
integrable, since it has finitely many finite values. Monotonicity
therefore bounds every lower simple integral by \(\int f\,d\mu\)
and every upper simple integral from below by that number.

For \(n\geq1\), on \(E\) define
\(a_n=2^{-n}\lfloor2^n f\rfloor\), where \(\lfloor x\rfloor\)
is the greatest integer at most \(x\), and set \(a_n=0\) outside
\(E\). Put \(b_n=a_n+2^{-n}\mathbf1_E\).
Boundedness of \(f\) gives a finite range for \(a_n\); its level
sets are intersections of \(E\) with inverse images of half-open
real intervals, so it is simple. The same holds for \(b_n\).
They satisfy \(a_n\leq f\leq b_n\), and
\[
 \int b_n\,d\mu-\int a_n\,d\mu=2^{-n}\mu(E)\longrightarrow0.
\]
Let \(L\) be the supremum of lower simple integrals and \(H\)
the infimum of upper simple integrals in the statement. The bounds
already proved give \(L\leq\int f\,d\mu\leq H\), while these
particular approximations give \(0\leq H-L\leq2^{-n}\mu(E)\).
All terms are finite: the approximations bound both \(L,H\) by
finite numbers. Letting \(n\) increase proves their equality with
the integral, including \(M=0\) or \(\mu(E)=0\).

For the sigma-finite assertion, the functions
\(h_n=\min(f,n)\mathbf1_{E_n}\) are nonnegative measurable,
bounded, and have positive regions contained in \(E_n\). They
increase pointwise to \(f\): each outcome lies in all sufficiently
large \(E_n\), and its truncated value tends to its value in
\([0,\infty]\). Monotone convergence proves the displayed limit.
Every bounded nonnegative \(h\leq f\) with finite-measure positive
region has integral at most \(\int f\,d\mu\), whereas the
admissible \(h_n\) have integrals increasing to it. Thus the
supremum over that class is exactly the same value.

Durrett's first stage uses signed simple functions on finite-measure
pieces, for which the signed integral already proved is the same finite
weighted sum. The first assertion gives his bounded stage, and the
second gives his nonnegative stage on a sigma-finite exhaustion. Both
constructions define the integrable signed stage by the same positive
and negative part subtraction. This establishes the stated agreement
without assuming it in any earlier proof.
\end{proof}
\end{annexitem}

\begin{annexitem}{annex:p15.dominated-bounded-convergence}{thm:p15-dominated-bounded-convergence}{Dominated convergence in L1 and bounded convergence in measure}
\begin{proof}
For the first assertion, remove one measurable null set outside which
all the countably many bounds and convergence hold. The set where
\(g=\infty\) is also null, because \(g\) is integrable. Set
\(f_n,f,g\) to zero on their union. These replacements are measurable
and do not change integrals or \(L^1\) distances. We may therefore
use finite real values everywhere, with \(|f_n|\leq g\), and
take the pointwise limit to get \(|f|\leq g\). All these functions
are integrable by monotonicity.

Put \(d_n=|f_n-f|\). Then \(0\leq d_n\leq2g\),
\(d_n\to0\), and \(2g-d_n\geq0\). Fatou's lemma gives
\[
 2\int g\,d\mu
 \leq\liminf_n\int(2g-d_n)\,d\mu
 =2\int g\,d\mu-\limsup_n\int d_n\,d\mu.
\]
All integrals subtracted here are finite and
\(0\leq\int d_n\,d\mu\leq2\int g\,d\mu\).
Subtracting the finite common term proves
\(\limsup_n\int d_n\,d\mu\leq0\), so the nonnegative sequence
converges to zero. The absolute-value inequality gives
\(\left|\int f_n\,d\mu-\int f\,d\mu\right|\leq\int d_n\,d\mu\),
proving convergence of the integrals too.

For the second assertion, first verify the properties of \(f\)
that were not initially assumed. For every \(\varepsilon>0\),
outside the null exceptions to the bounds on \(f_n\),
\[
 \{|f|>M+\varepsilon\}\subseteq\{|f-f_n|>\varepsilon\}.
\]
Taking measures and using convergence in measure shows that the set
on the left is null. Its union over \(\varepsilon=1/k\) gives
\(|f|\leq M\) a.e. Likewise
\(E^c\cap\{|f|>\varepsilon\}\subseteq\{|f-f_n|>\varepsilon\}\)
outside a null set, so \(f=0\) on \(E^c\) a.e. The bounds imply
integrability of \(f\) and of every \(f_n\).

Remove the countable union of their null exceptions; all integrals
and measures of disagreement sets are unchanged. Then for any
\(\delta>0\), split \(E\) into the sets where
\(|f_n-f|\leq\delta\) and \(|f_n-f|>\delta\). The integral
laws give
\[
 \|f_n-f\|_1\leq\delta\mu(E)
            +2M\mu(\{|f_n-f|>\delta\}).
\]
For fixed \(\delta\), the second term tends to zero. The first
can be made arbitrarily small because \(\mu(E)<\infty\).
This proves \(L^1\) convergence, and the same absolute-value bound
proves convergence of integrals. Finally, on a finite measure space a
finite constant is integrable, giving the stated bounded a.e. case
directly from the first assertion.
\end{proof}
\end{annexitem}

\begin{annexitem}{annex:p15.holder-integral}{thm:p15-holder-integral}{H\"older's inequality with its endpoint and zero cases}
\begin{proof}
For finite \(p\), \(\|f\|_p=0\) implies \(f=0\) a.e. by
the nonnegative zero-integral criterion. For the essential supremum,
if \(C=\|f\|_\infty<\infty\), the defining infimum provides
for each \(k\) a bound \(M_k<C+1/k\) with
\(\mu(\{|f|>M_k\})=0\). Therefore each
\(\{|f|>C+1/k\}\) is null, and their union is
\(\{|f|>C\}\). Thus \(|f|\leq C\) a.e., including when
\(C=0\). If either norm in the theorem is zero, one factor is
zero a.e.; since both functions are real valued, their product is
zero a.e. Its integral and the adopted right side are both zero.
If neither norm is zero and at least one is infinite, the right side
is infinity and the inequality holds in the extended order.

It remains to consider finite positive norms. Suppose first
\(1<p,q<\infty\). Then \(q=p/(p-1)\). For \(x,y\geq0\),
Young's inequality is
\[
 xy\leq x^p/p+y^q/q.
\]
For \(y=0\) it is immediate. For \(y>0\), the function
\(r(x)=x^p/p-xy+y^q/q\) has derivative \(x^{p-1}-y\),
negative below \(x_0=y^{1/(p-1)}\) and positive above it.
Its minimum on \([0,\infty)\) is therefore at \(x_0\).
Substitution, using \(x_0^p=y^q\) and \(x_0y=y^q\), gives
\(r(x_0)=y^q(1/p+1/q-1)=0\), proving the inequality.
Apply it pointwise to
\(x=|f|/\|f\|_p\), \(y=|g|/\|g\|_q\).
Nonnegative scaling, additivity and monotonicity give
\[
 \frac{\int|fg|\,d\mu}{\|f\|_p\|g\|_q}
  \leq\frac1p\frac{\int|f|^p\,d\mu}{\|f\|_p^p}
      +\frac1q\frac{\int|g|^q\,d\mu}{\|g\|_q^q}
  =\frac1p+\frac1q=1.
\]
The finite bound proves integrability and the desired inequality.

If \(p=1,q=\infty\), the essential-bound fact already proved
gives \(|fg|\leq\|g\|_\infty|f|\) a.e. Integrating proves
the result. Interchanging \(f,g\) handles \(p=\infty,q=1\).
These are all possibilities for the endpoints. The choice \(p=q=2\)
is the stated Cauchy--Schwarz case.
\end{proof}
\end{annexitem}

\begin{annexitem}{annex:p15.jensen-integral}{thm:p15-jensen-integral}{A supporting affine function and Jensen's inequality}
\begin{proof}
For \(x<y<z\), express \(y\) as the convex combination of \(x,z\)
with weights \((z-y)/(z-x)\) and \((y-x)/(z-x)\). Rearranging
the defining convex inequality gives the two secant comparisons
\[
 \frac{\varphi(y)-\varphi(x)}{y-x}
 \leq\frac{\varphi(z)-\varphi(x)}{z-x}
 \leq\frac{\varphi(z)-\varphi(y)}{z-y}.
\]
In particular, all slopes between \(c\) and points within distance
one of \(c\) lie between the finite slopes on \([c-1,c]\)
and \([c,c+1]\). Their absolute values are bounded by the larger
absolute value of those two fixed slopes. Thus
\(|\varphi(x)-\varphi(c)|\leq M_c|x-c|\) for \(|x-c|<1\),
with a finite \(M_c\). This proves continuity at each \(c\),
hence Borel measurability and measurability of \(\varphi(f)\).

Fix \(c\in\mathbb R\) and let
\[
 a=\sup_{x<c}\frac{\varphi(c)-\varphi(x)}{c-x}.
\]
The set is nonempty and bounded above by the slope on \([c,c+1]\),
and includes the finite slope on \([c-1,c]\); hence \(a\) is
a finite real number. Every left slope is at most \(a\), while
every right slope is at least each left slope and therefore at least
\(a\), by the secant comparisons. Rearranging both inequalities
gives, for every real \(x\),
\[
 \varphi(x)\geq\varphi(c)+a(x-c).
\]
Equality holds at \(x=c\). This constructs the supporting affine
function without a differentiability assumption on \(\varphi\).

Now take \(c=\int f\,d\mu\), which is finite by integrability.
The affine function of \(f\) on the right is integrable because
\(\mu(\Omega)=1\). Integrate the supporting inequality and use
signed monotonicity and linearity:
\[
 \int\varphi(f)\,d\mu
  \geq\varphi(c)\mu(\Omega)+a\left(\int f\,d\mu-c\mu(\Omega)\right)
  =\varphi(c).
\]
This is Jensen's inequality with exactly the stated integrability
and normalization hypotheses.
\end{proof}
\end{annexitem}

\begin{annexitem}{annex:p15.expectation-law-integral}{thm:p15-expectation-law-integral}{Expectation and integration on the value space}
\begin{proof}
Expectation is the same integral with \(\mu=\mathbb P\), so
the preceding theorems apply with their respective hypotheses unchanged.
The indicator identity is the simple integral identity. Constants
are integrable and have expectation equal to their value because
\(\mathbb P(\Omega)=1\). This gives the affine formula for
integrable \(X\) by signed linearity.

Suppose \(X\geq0\) and \(\mathbb E[X]=\infty\). For \(a>0\),
put \(H=aX+b\). Its negative part is bounded by \(b^-=\max(-b,0)\),
so it has finite expectation. Pointwise, in nonnegative arithmetic,
\(H^++b^-=aX+H^-+b^+\). Integrating proves
\(\mathbb E[H^+]+b^-=\infty\), whence
\(\mathbb E[H^+]=\infty\) and \(\mathbb E[H]=\infty\).
For \(a<0\), apply this argument to \((-a)X-b\) and then
exchange positive and negative parts to negate it. For \(a=0\),
the random variable is the constant \(b\), using the specified
zero-times-infinity convention if needed. This covers every branch
without subtracting two infinite expectations.

For the law formula, first take \(h=\mathbf1_B\) with
\(B\in\mathcal S\). Then
\[
 \mathbb E[\mathbf1_B(X)]=\mathbb P(X\in B)
     =\mathbb P_X(B)=\int_S\mathbf1_B\,d\mathbb P_X.
\]
Nonnegative finite linear combinations prove it for nonnegative simple
\(h\). For a nonnegative measurable \(h\), choose simple
\(h_n\uparrow h\) on the value space. Their composites increase
to \(h(X)\), and monotone convergence on each space gives
\[
 \mathbb E[h(X)]=\lim_n\mathbb E[h_n(X)]
   =\lim_n\int_S h_n\,d\mathbb P_X=\int_S h\,d\mathbb P_X.
\]
This includes an infinite integral. Applying this equality to \(|h|\)
proves equivalence of the two integrability conditions. When they
hold, apply the nonnegative equality to \(h^+\) and \(h^-\)
and subtract the two finite numbers. This proves the signed case.

For an atomic source space, first note that for any \(c\in[0,\infty]\)
and measurable \(A\),
\(\int c\mathbf1_A\,d\mathbb P=c\mathbb P(A)\).
Finite \(c\) uses the simple formula; for \(c=\infty\),
apply monotone convergence to \(n\mathbf1_A\). Its limiting integral
is zero if \(\mathbb P(A)=0\) and infinity otherwise, exactly
the convention in the formula. Now decompose a nonnegative function
into its countably many singleton-supported pieces. The nonnegative
series theorem gives
\(\mathbb E[X]=\sum_\omega X(\omega)p(\omega)\).
For an integrable real function, apply this result to its absolute
value and its two parts. Absolute convergence and finite subtraction
give the signed sum as stated.
\end{proof}
\end{annexitem}

\begin{annexitem}{annex:p15.expectation-bounds}{thm:p15-expectation-bounds}{Nonnegative probability bounds and variance identities}
\begin{proof}
Since \(A\) is nonempty and \(\varphi\) is finite valued and
nonnegative, \(c_A\) is a finite nonnegative number. Measurability
of \(\varphi\) and \(X\) makes \(\varphi(X)\) measurable.
The pointwise inequalities
\[
 c_A\mathbf1_{\{X\in A\}}
 \leq\varphi(X)\mathbf1_{\{X\in A\}}\leq\varphi(X)
\]
give the general bound by nonnegative monotonicity and the indicator
formula. An empty \(A\) has an empty event and zero restricted
integral, as separately stipulated. For nonnegative \(Y\), including
extended values, the bound \(r\mathbf1_{\{Y\geq r\}}\leq Y\)
and division by the positive finite \(r\) give Markov's inequality.

If \(\mathbb E[X^2]<\infty\), apply Cauchy--Schwarz to \(X\)
and the constant one. Its second norm is one, so
\(\mathbb E[|X|]\leq(\mathbb E[X^2])^{1/2}<\infty\).
The mean \(m\) is therefore defined and finite. The inequality
\((X-m)^2\leq2X^2+2m^2\) proves integrability of its square.
Expansion and signed linearity, now with every term integrable, give
\[
 \operatorname{Var}(X)=\mathbb E[X^2]-2m\mathbb E[X]+m^2
                     =\mathbb E[X^2]-m^2\geq0.
\]
The last inequality also follows directly from its defining nonnegative
square. Similarly \((aX+b)^2\leq2a^2X^2+2b^2\) has finite
expectation. Its mean is \(am+b\), so its centered variable is
\(a(X-m)\). Squaring and using nonnegative scaling proves
\(\operatorname{Var}(aX+b)=a^2\operatorname{Var}(X)\), including
\(a=0\). Finally apply Markov to \(Y=(X-m)^2\) at threshold
\(r^2>0\). Its threshold event is \(\{|X-m|\geq r\}\),
which gives Chebyshev's variance bound.
\end{proof}
\end{annexitem}

\begin{annexitem}{annex:p15.moment-controlled-convergence}{thm:p15-moment-controlled-convergence}{Truncation under a stronger-moment bound}
\begin{proof}
Continuity of \(g\) and almost-sure convergence give
\(g(X_n)\to g(X)\) a.s. Fatou and the assumed uniform bound yield
\(\mathbb E[g(X)]\leq\liminf_n\mathbb E[g(X_n)]\leq K\).
Replace \(h\) temporarily by \(h-h(0)\). The tail ratio still
tends to zero, because \(g\to\infty\), and adding back the
constant will preserve the desired conclusion. Hence suppose \(h(0)=0\).

For all sufficiently large \(M\), \(g(x)>0\) when \(|x|\geq M\).
Set \(\varepsilon_M=\sup_{|x|>M}|h(x)|/g(x)\). The tail-ratio
hypothesis makes this finite for large \(M\), with
\(\varepsilon_M\to0\). Continuity makes \(h\) bounded on
\([-M,M]\): its locally bounded neighborhoods have a finite
subcover by the compact-interval fact proved in the measure-extension
chapter. Let \(C_M\) bound \(|h|\) there. For any real random
variable \(Y\) with \(\mathbb E[g(Y)]\leq K\),
\[
 \mathbb E[|h(Y)|]\leq C_M+\varepsilon_M\mathbb E[g(Y)]<\infty.
\]
This proves integrability for \(Y=X_n\) and \(Y=X\) before
any signed comparison of their expectations.

At most countably many real values \(x\) satisfy \(\mathbb P(X=x)>0\), by
Theorem~\ref{thm:p15-cdf-laws}. We may therefore choose arbitrarily
large \(M\) for which \(\mathbb P(|X|=M)=0\), also satisfying
the preceding positivity condition on \(g\). For such an \(M\),
write \(Y^{(M)}=Y\mathbf1_{\{|Y|\leq M\}}\). At an outcome
where \(X_n\to X\) and \(|X|\ne M\), the truncations converge:
if \(|X|<M\) they eventually equal the original variables, and
if \(|X|>M\) they eventually vanish. Thus
\(h(X_n^{(M)})\to h(X^{(M)})\) a.s. All these variables are
bounded in absolute value by \(C_M\), so bounded convergence on
the probability space gives convergence of their expectations.

Since \(h(0)=0\), the truncation error for any such \(Y\) satisfies
\[
 \left|\mathbb E[h(Y)]-\mathbb E[h(Y^{(M)})]\right|
 \leq\mathbb E[|h(Y)|;|Y|>M]
 \leq\varepsilon_M\mathbb E[g(Y)]\leq K\varepsilon_M.
\]
Apply this to \(X_n\) and \(X\), and insert their truncated
expectations between the original ones. The triangle inequality and
the convergence at fixed \(M\) give
\[
 \limsup_n\left|\mathbb E[h(X_n)]-\mathbb E[h(X)]\right|
                         \leq2K\varepsilon_M.
\]
Let \(M\) increase through the permissible values. The right side
tends to zero, proving the assertion. Adding back \(h(0)\) returns
the original \(h\). Finally \(g(x)=|x|^p\), \(p>1\), and
\(h(x)=x\) satisfy the assumptions because
\(|x|/|x|^p=|x|^{1-p}\to0\) away from zero at infinity.
\end{proof}
\end{annexitem}
